\documentclass[10pt,a4paper]{amsart}
\usepackage[margin=19mm]{geometry}
\usepackage{amsmath,amssymb,mathtools}
\usepackage[scr=rsfs,cal=euler]{mathalfa}
\usepackage{xfrac}
\usepackage[unicode,hidelinks,bookmarks=false]{hyperref}
\newcommand{\ie}{i.e.\ }
\newcommand{\cf}{cf.\ }
\newcommand{\resp}{respectively\ }
\newcommand{\Subsection}{\subsection}
\providecommand{\nobreakdash}{\nobreak\hbox{-}}
\setlength{\emergencystretch}{2em}
\multlinegap0pt
\usepackage{amssymb}
\newtheorem{lemma}{Lemma}
\newcommand{\dd}{\mathrm{d}}
\title{Next-order asymptotics for the volume of \\ Schatten balls}
\author{Mathias Sonnleitner}
\date{Source: arXiv:2512.04933v2 (CC BY 4.0)}
\begin{document}
\maketitle

\noindent\emph{Source and licence.} Next-order asymptotics for the volume of \\ Schatten balls. Version arXiv:2512.04933v2 (CC BY 4.0), distributed by arXiv under the Creative Commons Attribution 4.0 International licence, http://creativecommons.org/licenses/by/4.0/, as recorded in the arXiv licence metadata for this exact version and shown on the abstract page https://arxiv.org/abs/2512.04933v2. Full licence notice: \texttt{licenses/arxiv-2512.04933-CC-BY-4.0.txt}.\par\smallskip

\noindent\textit{Editorial scope.} Counted unit: the article's lemma lem:hoelder (source0.tex lines 319--324). The article's complete original proof of it is delivered with it (source0.tex lines 325--365, one \texttt{proof} environment of 41 lines). No mathematics is written, repaired or re-derived: the statement and the proof are the article's own lines, in the article's own notation, and no retained line is substituted or re-flowed.  No further source-local context is needed: every cross-reference of every retained line resolves inside the two delivered spans.  Nothing was omitted from any retained span, so the package carries no omission record.  No retained line contains a citation, so the wrapper carries no bibliography and no bibliography entry.  No retained line contains a figure environment, an image-inclusion command or drawing code, so no image is delivered and the package declares no asset dependency.  Editorial formatting only, in addition: an AMS \texttt{amsart} preamble that restates the article's own declarations and macros used by the retained lines, namely the environments \texttt{lemma} on separate counters, together with the standard packages those lines need.  The retained environments print in this document as Lemma 1 in order of appearance, whereas the article prints the same items with the numbering of its own full document.  The complete native source of the article as distributed in the arXiv e-print for this exact version is bundled unchanged as \texttt{sources/190364/source0.tex}.
\begin{lemma} \label{lem:hoelder}
	Let $p>1$. Then $f_p$ is Hölder continuous of order $\min\{p-1,\frac{1}{2}\}$ on $[-1,1]$. Moreover, it is Hölder continuous of order $\frac{1}{2}$ in neighborhoods of $-1$ and $1$ and there exist positive constants $c$ and $C$ such that 
\begin{equation} \label{eq:tail}
	c\sqrt{1-|x|}\le f_p(x) \le C\sqrt{1-|x|}\qquad \text{for}\qquad |x|\ge\frac{1}{2}.
\end{equation}
\end{lemma}

\begin{proof}
First, we verify the behavior of $f_p$ near the boundary points. By a change of variables, for $x\in (-1,1)\setminus\{0\}$, 
\begin{equation} \label{eq:change}
\frac{\pi}{p}f_p(x)
=\int_{|x|}^1 \frac{t^{p-1}}{\sqrt{t^2-x^2}}\dd t
=|x|^{p-1}\int_0^{\arccos |x|}\frac{1}{\cos^p \alpha}\dd \alpha.
\end{equation}
Combining \eqref{eq:change} with the asymptotics $\arccos |x|=\sqrt{2(1-|x|)}+o(1-|x|)$ as $|x|\to 1$ proves \eqref{eq:tail} which in turn implies that $f_p$ is Hölder continuous of order $\frac{1}{2}$ at the boundary points $\pm 1$.

Next, we show that $f_p$ is in fact continuously differentiable at $x\in (-1,1)\setminus \{0\}$. For symmetry reasons it is enough to consider $x\in (0,1)$. Let $0<|h|<\min\{x,1-x\}$ and write 
\begin{align*}
\frac{\pi}{p}\big(f_p(x+h)-f_p(x)\big)	
=&\big((x+h)^{p-1}-x^{p-1}\big)\int_0^{\arccos x}\frac{1}{\cos^p\alpha}\dd\alpha\\
&-(x+h)^{p-1}\int_{\arccos(x+h)}^{\arccos x}\frac{1}{\cos^p \alpha}\dd \alpha.
\end{align*}
Using $\arccos(x+h)=\arccos x -\frac{h}{\sqrt{1-x^2}}+o(h)$ and letting $h\to 0$ shows that $f_p$ is continuously differentiable at $x$ with derivative
\[
f_p'(x)
=\frac{p-1}{x}f_p(x)-\frac{p}{\pi}\frac{1}{x\sqrt{1-x^2}}.
\]

It remains to check Hölder continuity at zero. Let $h\in (0,1)$ and write
\begin{align*}
\frac{\pi}{p}\big(f_p(h)-f_p(0)\big)
&=\int_h^1 \frac{t^{p-1}}{\sqrt{t^2-h^2}}\dd t - \int_0^1 t^{p-2}\dd t\\
&=\int_h^1 t^{p-1}\Big(\frac{1}{\sqrt{t^2-h^2}}-\frac{1}{t}\Big)\dd t -\frac{h^{p-1}}{p-1}.
\end{align*}
The latter integral can be estimated by
\begin{align*}
\int_h^1 t^{p-1}\Big(\frac{1}{\sqrt{t^2-h^2}}-\frac{1}{t}\Big)\dd t
&\le h^2\int_h^1 \frac{t^{p-3}}{\sqrt{t^2-h^2}}\dd t\\
&= h^{p-1}\int_1^{1/h} \frac{u^{p-3}}{\sqrt{u^2-1}}\dd u.
\end{align*}
In particular, if $1< p<2$, then
\begin{equation} \label{eq:hoelder}
0\le \frac{\pi}{p}\frac{f_p(h)-f_p(0)}{h^{p-1}}
\le \int_1^{\infty} \frac{u^{p-3}}{\sqrt{u^2-1}}\dd u
<\infty.
\end{equation}
Thus, in this case we obtain that $f_p$ is Hölder continuous of order $p-1$ at $0$. For $p\ge 2$, a similar argument shows that $f_p$ is continuously differentiable at $0$ with $f_p'(0)=0$. In conclusion, if $p>1$, the function $f_p$ is Hölder continuous of order $\min\{p-1,\frac{1}{2}\}$ on $[-1,1]$. This completes the proof.
\end{proof}

\end{document}
